\chapter{\IfLanguageName{dutch}{Stand van zaken}{State of the art}}%
\label{ch:stand-van-zaken}

Dit hoofdstuk situeert de bachelorproef binnen het bestaande wetenschappelijke en technologische kader rond AI-gedreven conversatiesystemen, observability en AI-specifieke beveiliging. In het vorige hoofdstuk werd de probleemstelling geformuleerd, waarbij werd vastgesteld dat AI-chatbots steeds vaker worden geïntegreerd in enterprise-omgevingen, maar dat monitoring en beveiliging vaak los van elkaar worden geïmplementeerd. \\

De literatuurstudie onderzoekt daarom de onderliggende concepten, architecturen en beveiligingsuitdagingen die relevant zijn voor het ontwerpen van een geïntegreerde oplossing. Eerst wordt de werking en evolutie van AI-chatbots besproken, gevolgd door de architecturale kenmerken van Large Language Models in enterprise-context. Vervolgens worden de belangrijkste securityrisico’s en AI-firewall\-mechanismen behandeld. Daarna wordt ingezoomd op observability, logging en eventcorrelatie binnen gedistribueerde systemen. Tot slot worden privacy- en ethische aspecten besproken en wordt het voorgestelde framework gepositioneerd ten opzichte van bestaande oplossingen.\\

Dit hoofdstuk vormt de conceptuele basis voor de methodologie in het volgende hoofdstuk, waarin wordt toegelicht hoe deze inzichten worden toegepast in het ontwerp en de implementatie van de proof-of-concept.\\

\section{Inleiding tot AI-gedreven Conversatiesystemen}

AI-gedreven conversatiesystemen, vaak aangeduid als chatbots of virtuele assistenten, zijn softwaretoepassingen die natuurlijke taal gebruiken om interactie met gebruikers mogelijk te maken. Waar vroege chatbots voornamelijk werkten op basis van vooraf gedefinieerde regels en beslissingsbomen, maken moderne systemen gebruik van machine learning en, meer recent, Large Language Models (LLM’s). Deze evolutie heeft geleid tot een significante toename in flexibiliteit, schaalbaarheid en contextbegrip binnen menselijke-computerinteractie \autocite{adamopoulou2020chatbots}.\\

De opkomst van generatieve AI heeft het toepassingsgebied van chatbots verder uitgebreid. In tegenstelling tot traditionele Natural Language Processing (NLP)-systemen, die vaak beperkt zijn tot intentieherkenning en vooraf geprogrammeerde antwoorden, kunnen LLM-gebaseerde systemen autonoom tekst genereren op basis van context en probabilistische taalmodellen \autocite{brown2020language}. Hierdoor zijn ze in staat om complexere en meer natuurlijke dialogen te voeren.\\

\subsection{Evolutie van Chatbots}

De eerste generatie chatbots, zoals ELIZA in de jaren 60, functioneerden op basis van patroonherkenning en eenvoudige teksttransformaties \autocite{weizenbaum1966eliza}. Deze systemen simuleerden conversatie zonder daadwerkelijk begrip van taalstructuur of semantiek. In de daaropvolgende decennia werden rule-based systemen uitgebreid met decision trees en keyword matching, wat leidde tot meer gestructureerde maar nog steeds beperkte interacties.\\

% semantiek = betekenisleer

Met de opkomst van machine learning en statistische NLP-technieken verschoof de focus naar intentieclassificatie en entiteitsextractie. Platforms zoals dialog management frameworks maakten het mogelijk om schaalbare klantenservicebots te ontwikkelen. Toch bleven deze systemen sterk afhankelijk van gestructureerde trainingsdata en vooraf gedefinieerde intenties \autocite{jurafsky2023speech}.\\

De introductie van transformer-gebaseerde modellen betekende een fundamentele verschuiving. Het transformer-architectuurmodel, geïntroduceerd door \textcite{vaswani2017attention}, maakte het mogelijk om lange-afstandsafhankelijkheden in tekst efficiënt te modelleren. Hieruit ontstonden Large Language Models zoals GPT, die getraind worden op enorme tekstcorpora en in staat zijn coherente, contextbewuste tekst te genereren \autocite{brown2020language}. Dit zijn grote, gestructureerde verzameling van geschreven of gesproken teksten in een bepaalde taal. Deze modellen vormen vandaag de basis van moderne AI-chatbots.\\

\subsection{Toepassingen in Enterprise-omgevingen}

Binnen enterprise-omgevingen worden AI-chatbots steeds vaker ingezet voor uiteenlopende toepassingen, waaronder klantenondersteuning, interne helpdesks, kennisbeheer en procesautomatisering. De schaalbaarheid en 24/7-beschikbaarheid maken deze systemen economisch aantrekkelijk \autocite{dwivedi2023opinion}.\\

Tegelijkertijd brengt de integratie van generatieve AI binnen bedrijfsomgevingen nieuwe uitdagingen met zich mee. In tegenstelling tot klassieke systemen genereren LLM’s dynamische output, wat betekent dat antwoorden niet volledig voorspelbaar zijn. Dit verhoogt het risico op foutieve informatie, ongepaste inhoud of het lekken van gevoelige data \autocite{bommasani2021opportunities}.\\

Bovendien worden AI-chatbots vaak geïntegreerd via API’s in bredere IT-architecturen, waarbij zij toegang krijgen tot interne databronnen. Hierdoor verschuift het risicoprofiel van enkel functionele foutmarges naar potentiële security- en complianceproblemen. \\

\subsection{Noodzaak van Monitoring en Beveiliging}

De probabilistische aard van generatieve AI maakt traditionele kwaliteitscontrolemechanismen onvoldoende. Waar klassieke software deterministisch gedrag vertoont, kunnen LLM’s verschillende antwoorden genereren op identieke input. Dit bemoeilijkt validatie, auditing en incidentanalyse. \\

Onderzoekers wijzen erop dat observability en security in AI-systemen vaak als afzonderlijke componenten worden benaderd, terwijl een geïntegreerde aanpak noodzakelijk is om anomalieën effectief te detecteren en te analyseren \autocite{bommasani2021opportunities}. Logging van conversaties, traceerbaarheid van interacties en correlatie van security-events zijn cruciale elementen om verantwoord gebruik van AI-systemen te waarborgen. \\

De toenemende afhankelijkheid van AI binnen kritieke bedrijfsprocessen onderstreept daarom de noodzaak van geïntegreerde monitoring- en beveiligingsmechanismen. Dit vormt de basis voor het verdere onderzoek in deze bachelorproef.\\


\section{Architectuur van Large Language Models in Enterprise-omgevingen}

Large Language Models (LLM’s) vormen de technologische kern van moderne AI-gedreven conversatiesystemen. Deze modellen zijn gebaseerd op deep learning-architecturen en worden getraind op grootschalige tekstcorpora om probabilistische taalpatronen te leren. Binnen enterprise-omgevingen worden LLM’s doorgaans niet als standalone systemen gebruikt, maar geïntegreerd als component binnen een bredere IT-architectuur.

Dit hoofdstuk bespreekt de onderliggende architectuurprincipes van LLM’s, hun integratiemodellen binnen bedrijfsomgevingen en de gevolgen voor monitoring en beveiliging.

\subsection{Basisprincipes van Large Language Models}

LLM’s zijn neurale netwerken met miljarden parameters die taal modelleren op basis van waarschijnlijkheidsverdelingen. In plaats van vaste regels toe te passen, voorspellen zij het volgende token in een sequentie op basis van context \autocite{brown2020language}. Hierdoor kunnen zij coherente en contextbewuste tekst genereren.

De schaal van deze modellen is een bepalende factor voor hun performantie. \textcite{kaplan2020scaling} tonen aan dat modelprestaties systematisch verbeteren naarmate het aantal parameters, trainingsdata en rekencapaciteit toenemen. Dit schaalprincipe verklaart de sterke vooruitgang in generatieve AI-systemen.

\subsection{Transformer-architectuur}

De meeste hedendaagse LLM’s zijn gebaseerd op de transformer-architectuur, geïntroduceerd door Vaswani et al. \cite{vaswani2017attention}. Het kernmechanisme van deze architectuur is self-attention, waarmee het model relevante woorden in een zin kan wegen ten opzichte van elkaar, ongeacht hun positie in de tekst.

In tegenstelling tot recurrente neurale netwerken maakt de transformer parallelle verwerking mogelijk, wat schaalbaarheid en efficiëntie aanzienlijk verhoogt. Dit maakt het trainen van zeer grote modellen haalbaar op gedistribueerde infrastructuren.

Voor enterprise-toepassingen betekent dit dat de intelligentie voornamelijk in het model zelf zit, terwijl de applicatielaag zich richt op interfacing, validatie, logging en beveiliging.

\subsection{Integratie via API-gebaseerde Architecturen}

Binnen bedrijfsomgevingen worden LLM’s vaak geïntegreerd via API’s. De applicatie dient hierbij als tussenlaag tussen gebruiker en model. De typische dataflow verloopt als volgt:

\begin{enumerate}
    \item Gebruiker stuurt een prompt naar de applicatie.
    \item De applicatie valideert de input.
    \item De prompt wordt doorgestuurd naar het LLM via een API-call.
    \item Het model genereert een response.
    \item De applicatie verwerkt, filtert en logt de output.
\end{enumerate}

Dit architectuurmodel verduidelijkt dat beveiliging en observability niet in het model zelf zitten, maar in de omliggende infrastructuur. Hierdoor ontstaat een duidelijke noodzaak om input- en outputstromen systematisch te monitoren.

\subsection{SaaS versus Self-hosted Implementaties}

Organisaties kunnen kiezen tussen cloudgebaseerde (SaaS) LLM-oplossingen of self-hosted modellen. SaaS-oplossingen bieden schaalbaarheid en eenvoud in beheer, maar brengen afhankelijkheid van externe providers met zich mee. Self-hosted modellen bieden meer controle over data en configuratie, maar vereisen aanzienlijke infrastructuur en expertise \cite{bommasani2021opportunities}.

De keuze tussen beide modellen heeft directe implicaties voor:

\begin{itemize}
    \item Databeheer en compliance
    \item Loggingmogelijkheden
    \item Beveiligingscontroles
    \item Observability-integratie
\end{itemize}

\subsection{Dataflow binnen AI-Chatbots}

De integratie van LLM’s in enterprise-systemen gebeurt zelden zonder bijkomende componenten. Vaak worden retrieval-mechanismen toegevoegd (Retrieval-Augmented Generation, RAG), waarbij externe databronnen worden geraadpleegd om context te verrijken \cite{lewis2020rag}. 

Dit introduceert extra complexiteit in de dataflow:
\begin{itemize}
    \item Query naar kennisdatabase
    \item Verrijking van prompt
    \item Modelinferentie
    \item Outputvalidatie
\end{itemize}

Elke stap vormt een potentieel controlepunt voor logging en beveiliging.

\subsection{Architecturale Integratie van Monitoring en Firewalling}

Aangezien LLM’s probabilistisch gedrag vertonen, kunnen ongewenste outputs of manipulaties niet volledig binnen het model zelf worden uitgesloten. Daarom wordt in recente literatuur gepleit voor een ``defence-in-depth''-benadering rond generatieve AI-systemen \cite{bommasani2021opportunities}.

Dit betekent dat monitoring- en beveiligingslagen worden toegevoegd rond het model:

\begin{itemize}
    \item Input filtering (detectie van prompt injection)
    \item Output filtering (toxicity, data leakage)
    \item Logging en traceerbaarheid
    \item Event correlatie
\end{itemize}

Deze architecturale positionering vormt de basis voor de proof-of-concept dat in deze bachelorproef wordt ontwikkeld.


\section{Securityrisico’s bij AI-Chatbots}
\subsection{Prompt Injection}
\subsection{Jailbreaking van AI-modellen}
\subsection{Data Exfiltration via Modeloutput}
\subsection{Model Misuse en Abuse}
\subsection{Beperkingen van Traditionele Beveiligingsmechanismen}

\section{AI-Firewalls en Prompt Filtering}
\subsection{Definitie en Werking van AI-Firewalls}
\subsection{Input Filtering en Output Filtering}
\subsection{Detectie van Toxiciteit en Ongepaste Inhoud}
\subsection{Anomaliedetectie in Conversatiepatronen}
\subsection{Verschillen met Klassieke Web Application Firewalls}

\section{Observability in Gedistribueerde Systemen}
\subsection{Definitie van Observability}
\subsection{De Drie Pijlers: Logs, Metrics en Traces}
\subsection{Distributed Tracing in Microservices-architecturen}
\subsection{Monitoring versus Observability}
\subsection{Uitdagingen bij AI-gedreven Systemen}

\section{Logging, Traceerbaarheid en Monitoring van AI-Systemen}
\subsection{Structurering van Conversatielogs}
\subsection{User- en Session-identificatie}
\subsection{Tijdstempels en Eventregistratie}
\subsection{Analyse van AI-output (Toxiciteit en Anomalieën)}
\subsection{Dashboarding en Visualisatie}
\subsection{Correlatie tussen Monitoringdata en Security-events}

\section{Correlatie tussen Security-events en Observabilitydata}
\subsection{Principes van Event Correlatie}
\subsection{Integratie van Threat Intelligence}
\subsection{SIEM-benaderingen in AI-omgevingen}
\subsection{Detectie van Afwijkende Interactiepatronen}
\subsection{Incidentrespons binnen AI-systemen}

\section{Privacy- en Ethiekvraagstukken bij AI Monitoring}
\subsection{GDPR en Dataminimalisatie}
\subsection{Privacy-by-Design in AI-systemen}
\subsection{Bewaartermijnen en Logretentie}
\subsection{Transparantie en Gebruikersrechten}
\subsection{Ethiek en Bias in AI-monitoring}

\section{Bestaande Oplossingen en Technologische Positionering}
\subsection{AI Observability Platforms}
\subsection{AI Security en Firewalloplossingen}
\subsection{Enterprise Monitoring Frameworks}
\subsection{Vergelijkende Analyse van Bestaande Benaderingen}
\subsection{Positionering van het Voorgestelde Framework}